\section{稳健性检验与讨论}

\subsection{服务半径敏感性检验}

\subsubsection{15 分钟阈值调整}

将长者食堂和社区服务站阈值调整为 10 分钟、20 分钟。

\subsubsection{30 分钟阈值调整}

将日间照料服务阈值调整为 20 分钟、40 分钟。

\subsubsection{结果稳定性分析}

说明高错配区域是否保持稳定。

\subsection{代际共享修正系数敏感性检验}

\subsubsection{不同共享修正系数取值设定}

检验：

\begin{equation}
\alpha = 0.1,\ 0.2,\ 0.3
\end{equation}

\subsubsection{高错配区域变化}

比较不同 $\alpha$ 下的高错配网格是否变化明显。

\subsubsection{稳健性结论}

若主要高错配区域基本一致，说明模型稳定。

\subsection{模型对比检验}

\subsubsection{LightGBM 与传统模型对比}

可与随机森林、逻辑回归或 XGBoost 简单对比。

\subsubsection{GraphSAGE 与非图模型对比}

说明加入空间邻接信息后，识别效果是否提升。

\subsubsection{检验结论}

这一节不需要长篇，只需说明主结论稳健。
